% !PRINTED 495
\weberpage{495}
%% [page furniture, not body text -- source-confirmed 2026-08-10]
%% \S\ 133. \hfill The resolvents of degrees 7 and 11.

%<<CONT
$A_\nu$ remains unchanged under the simultaneous interchange:
\[
\left( f(\omega)^{12}, \frac{64}{f(\omega)^{12}} \right), \quad
\left( i\sqrt{11}, -i\sqrt{11} \right),
\]
and for even $\nu$ it contains only the even powers of
$f(\omega)^{12}$, for odd $\nu$ only the odd ones.

Here too only powers of $f(\omega)$ occur in the denominator, and the
powers of $f(\omega)$ rise to the same height as those of
$f(\omega)^{-1}$.

In order to come to a clear view of the degrees of the functions
$A_\nu$, we consider the beginnings of the developments in ascending
powers of $q = e^{\pi i \omega}$.

The development of $f(\omega)$ begins with $q^{-\frac{1}{24}}$, that
of $v_\infty$ with $q^{-\frac{n}{24}}$, that of $v_c$ with
$q^{-\frac{1}{24n}} e^{-\frac{2\pi i c}{48}\, \frac{1}{n}}$; but in
the last exponential quantity, if we reduce $c$ to its least residue
$(\mathrm{mod}\ n)$, $\frac{1}{48}$ is also to be taken according to
the modulus $n$, and so for $n = 7, 11$
\[
\frac{1}{48} \equiv -1 \ (\mathrm{mod}\ 7), \qquad
\frac{1}{48} \equiv 3 \ (\mathrm{mod}\ 11)
\]
is to be set. If accordingly we now set
\[
\varrho = e^{\frac{2\pi i}{7}} \quad \text{or} \quad = e^{-\frac{6\pi i}{11}}
\]
then we obtain from (14), (15) as the beginnings of the development
for $n=7$:
\[
i\sqrt{7}\, w_\nu = q^{-\frac{1}{7}} \varrho^{3\nu} (\varrho - \varrho^3)(\varrho^2 - \varrho^6)(\varrho^4 - \varrho^5) + \cdots,
\]
for $n=11$:
\[
i\sqrt{11}\, w_\nu = q^{-\frac{5}{22}} \varrho^{5\nu} (\varrho - \varrho^2)(\varrho^3 - \varrho^6)(\varrho^4 - \varrho^8)(\varrho^5 - \varrho^{10})(\varrho^9 - \varrho^7) + \cdots
\]
But for $n = 7$:
\[
(\varrho - \varrho^3)(\varrho^2 - \varrho^6)(\varrho^4 - \varrho^5) = -\varrho - \varrho^2 - \varrho^4 + \varrho^3 + \varrho^5 + \varrho^6,
\]
and for $n = 11$:
\[
(\varrho - \varrho^2)(\varrho^3 - \varrho^6)(\varrho^4 - \varrho^8)(\varrho^5 - \varrho^{10})(\varrho^9 - \varrho^7) =
\]
\[
-\varrho - \varrho^3 - \varrho^4 - \varrho^5 - \varrho^9 +.\varrho^2 + \varrho^6 + \varrho^8 + \varrho^7 + \varrho^{10}. % [sic] stray period before $\varrho^2$ as printed
\]
